% Introduction to Part VIII. Unnumbered text: no \label except the \phantomsection target.
\phantomsection\label{intro:VIII}
The covariance $C_k$, the coarse Hessian, the determinant and the background map of the step depend on the background everywhere on the torus, and the polymer expansion of Part~X needs them as sums of local pieces. Two kinds of localization occur (Section~\ref{sec:interaction-local}, Summary). In a localization of \emph{radius type}, a piece that depends on a ball of radius $m$ is bounded by $e^{-\gamma m}$. In one of \emph{volume type}, a piece carries a fixed small factor for every cube of its dependence region, as \cite[(116)--(122)]{DimockI} requires. A ball of radius $m$ contains of order $m^4$ cubes, so radius type does not suffice. Part~VIII builds the Dirichlet operators, the interior estimates and the graded domains; the localization by telescoping that it also proves is of radius type, and Part~IX replaces it by volume type.

\paragraph{What this Part supplies to the main line.} The status paragraphs of Sections~\ref{sec:localization} and~\ref{sec:det-local} list it:
\begin{itemize}
\item Dirichlet versions on domains of the averaged operators $\boldsymbol\Gamma^X$, $\bH^X$, $\bC^X$, of the coarse Hessian and of the fluctuation operator, with uniform coercivity and decay, dependence on the background only near the domain, and comparison of domains with errors concentrated at the seam (Propositions~\ref{prop:dirichlet-averaged}, \ref{prop:dirichlet-hessian} and~\ref{prop:local-fluct}, Lemmas~\ref{lem:hessian-gamma} and~\ref{lem:coercive-general}). Propositions~\ref{prop:dirichlet-averaged} and~\ref{prop:local-fluct} are the inputs (W2) of the walk levels (L4) and (L5) of Section~\ref{sec:nested-walks};
\item interior pointwise bounds for the Dirichlet operators, Theorem~\ref{thm:ID}(a)--(c), used in Sections~\ref{sec:pointwise-weak}--\ref{sec:step-quantities}. Part~III does not give them, because its bootstrap is global (Remark~\ref{rem:ID}). Section~\ref{sec:interior} proves them by an absorption at all scales for a local elliptic system, obtained by treating two scalar potentials of the slice projection as unknowns (Lemma~\ref{lem:local-system}, Proposition~\ref{prop:system-interior});
\item the local determinant densities $\ell_Y(b)$, the background and graded activities on the domains $\mathcal U(P;\theta_r,\delta_1,\delta_0)$ (Definitions~\ref{def:bg-activity} and~\ref{def:graded}), and the quadratic bound: a normalized graded activity is bounded by twice its norm times the square of a quantity $\mathfrak s$ built from the curvature and the current of the real part and from the complex part of the configuration, whenever $\mathfrak s\le\frac12$ (Lemma~\ref{lem:quadratic-bound}). These, with Lemma~\ref{lem:complex-gauge} and Remarks~\ref{rem:later-scales}(b) and~\ref{rem:bg-to-graded}, are the setting of the activities of Parts~IX--XI.
\end{itemize}
Remark~\ref{rem:domain-constants} fixes the constants of the graded domains, which do not depend on the coupling, the step, the torus or the polymer. It refers to Parts~IX and~X throughout and is best read with Part~X.

\paragraph{Superseded, or off the main line.} The radius-type results are not used later: Theorem~\ref{thm:local-covariance} and Corollary~\ref{cor:large-changes} (Section~\ref{subsec:local-cov-expansion}), with Lemma~\ref{lem:telescoping} behind them, are superseded by the weakened walks of Section~\ref{sec:nested-walks}; Proposition~\ref{prop:det-local} and Corollary~\ref{cor:one-loop} by Proposition~\ref{prop:one-loop-polymer}; and Section~\ref{sec:interaction-local} by the weakened background map of Section~\ref{sec:weak-background}. Theorem~\ref{thm:ID}(d) and the identity of Lemma~\ref{lem:div-curvature} serve only Section~\ref{sec:interaction-local}; the main line uses from that lemma the bound on the operator $W_0$. Of Remark~\ref{rem:later-scales}, the main line uses part~(b).

\paragraph{Imports.} From Part~III: domains, the Dirichlet operators, locality and the comparison of inverses (Theorems~\ref{thm:reg-scalar}--\ref{thm:reg-locality}, Lemma~\ref{lem:compare-inverse}). From Parts~IV--VII: the averaging, the action, the contraction of Section~\ref{sec:var-solution}, Corollary~\ref{cor:regularity}, and the covariance and determinant of the step (Corollary~\ref{cor:covariance}, Proposition~\ref{prop:det-densities}). The clauses on complex perturbations, in Propositions~\ref{prop:dirichlet-averaged}(e) and~\ref{prop:dirichlet-hessian}(c) and Lemma~\ref{lem:coercive-general}, use Section~\ref{sec:complex-references} and Proposition~\ref{prop:hessian-general}, and Remark~\ref{rem:later-scales}(d) uses Lemma~\ref{lem:torus}; they are best read after those results.

\paragraph{Details for this Part.} The companion, chapter on Part~VIII, opens with the list of what the main line uses and a seam calculus for the comparison of domains. It expands the Dirichlet operators, the interior estimates and Theorem~\ref{thm:ID}, the quadratic bound and the domain constants, and marks the superseded results.
